\documentclass[tikz,border=3mm]{standalone}

%%%% Serifenloser Textstil (Calibri-ähnlich)
\usepackage[T1]{fontenc}
\usepackage[sfdefault,lf]{carlito}
\usepackage{lmodern}
\renewcommand*\familydefault{\sfdefault}

\usepackage{amsmath}
\usetikzlibrary{positioning,calc,arrows.meta,decorations.pathreplacing}

% ===== Parameter =====
\def\imgW{3.0}
\def\imgH{1.6}
\def\funW{2.65}

\pgfmathsetmacro{\boxW}{16.2}
\pgfmathsetmacro{\boxH}{8.8}

% Linienstärken
\def\imgFrameLW{1.0pt}
\def\boxLW{1.0pt}

\begin{document}
\begin{tikzpicture}[
  >={Latex[length=3mm]},
  every node/.style = {font=\tiny},
  imgframe/.style={
    draw=black,
    line width=\imgFrameLW,
    inner sep=0pt,
    minimum width=\imgW cm,
    minimum height=\imgH cm,
    fill=gray!10,
    align=center,
    rounded corners=2pt
  },
  bigbox/.style={
    draw=green!60!black,
    fill=green!60!black,
    fill opacity=0.16,
    rounded corners=8pt,
    line width=\boxLW,
    inner sep=8pt
  },
  fun/.style={
    draw=black,
    rounded corners=3pt,
    line width=0.75pt,
    fill=white,
    fill opacity=0.97,
    text opacity=1,
    minimum height=8.5mm,
    align=center,
    font=\tiny\bfseries,
    minimum width=\funW cm
  },
  model/.style={
    draw=black,
    rounded corners=8pt,
    line width=0.8pt,
    fill=blue!22,
    fill opacity=0.95,
    text opacity=1,
    minimum height=9.5mm,
    minimum width=3.4cm,
    align=center,
    font=\tiny\bfseries
  },
  emph/.style={
    draw=black,
    rounded corners=3pt,
    line width=0.8pt,
    fill=orange!18,
    fill opacity=0.95,
    text opacity=1,
    minimum height=8.5mm,
    align=center,
    font=\tiny\bfseries,
    minimum width=2.8cm
  },
  smooth/.style={
    draw=black,
    rounded corners=3pt,
    line width=0.8pt,
    fill=teal!16,
    fill opacity=0.95,
    text opacity=1,
    minimum height=8.5mm,
    align=center,
    font=\tiny\bfseries,
    minimum width=3.0cm
  },
  flow/.style={-Latex, line width=0.9pt},
  altflow/.style={-Latex, line width=0.8pt, dashed, color=black!75},
  feedback/.style={-Latex, line width=0.85pt, rounded corners=10pt},
  flowlabel/.style={midway, fill=white, inner sep=1pt, font=\tiny}
]

% =================================================
% LINKE SEITE: EINGÄNGE
% =================================================
\node[imgframe] (ref) at (0, 2.4) {Referenz\\Keypoints};
\node[above=2pt of ref, anchor=south, font=\scriptsize] {Soll-Merkmale};

\node[imgframe] (cam) at (0,-2.4) {Kamera /\\aktuelles Bild};
\node[below=2pt of cam, anchor=north, font=\scriptsize] {Ist-Beobachtung};

% =================================================
% MITTE: HAUPTBOX
% =================================================
\node[bigbox, anchor=west, minimum width=\boxW cm, minimum height=\boxH cm] (ctrlBox)
  at (3.0,-0.15) {};

\node[font=\bfseries, anchor=north, text=black]
  at ([yshift=-4pt]ctrlBox.north) {Regelung und Merkmalsverarbeitung};

% Bereichstitel
\node[font=\scriptsize\bfseries, text=black!70, anchor=north west]
  at ([xshift=0.9cm,yshift=-0.5cm]ctrlBox.north west) {Merkmalsverarbeitung};

\node[font=\scriptsize\bfseries, text=black!70, anchor=north east]
  at ([xshift=-0.9cm,yshift=-0.5cm]ctrlBox.north east) {Regelung};

% leichte Trennlinie innerhalb der Hauptbox
\draw[black!25, line width=0.5pt]
  ([xshift=0.9cm,yshift=-1.05cm]ctrlBox.north west) --
  ([xshift=-0.9cm,yshift=-1.05cm]ctrlBox.north east);

% -------------------------------------------------
% Obere Reihe: Wahrnehmung und Fehler
% -------------------------------------------------
\node[fun]  (detect) at (5.3, 2.15) {Keypoint\\Detection};
\node[fun]  (match)  at (8.3, 2.15) {Feature\\Matching};
\node[fun]  (switch) at (11.3,2.15) {Confidence\\\& Switching};
\node[emph] (error)  at (14.4,2.15) {Bildfehler};

% -------------------------------------------------
% Untere Reihe: Vorhersage und Regelung
% -------------------------------------------------
\node[model]  (pred)   at (11.3,0.05) {Modellbasierte\\Keypoint-Vorhersage};
\node[fun]    (ibvs)   at (14.4,0.05) {IBVS-\\Regler};
\node[smooth] (smooth) at (17.6,0.05) {Glättung /\\Begrenzung};
\node[fun]    (twist)  at (20.8,0.05) {Twist-\\Transformation};

% =================================================
% RECHTE SEITE: AUSGANG / STRECKE
% =================================================
\node[imgframe, minimum width=3.2cm, minimum height=1.65cm] (robot)
  at (24.5, 0.95) {Roboterarm};
\node[above=2pt of robot, anchor=south, font=\scriptsize] {Ausgang};

\node[imgframe, minimum width=3.2cm, minimum height=1.65cm] (scene)
  at (24.5,-2.1) {Szene / Objekt};
\node[below=2pt of scene, anchor=north, font=\scriptsize] {Rückkopplung};

% =================================================
% HAUPTFLUSS
% =================================================
\draw[flow] (cam.east) -- (detect.west);
\draw[flow] (detect.east) -- (match.west);
\draw[flow] (match.east) -- (switch.west);
\draw[flow] (switch.east) -- (error.west);

\draw[flow] (error.south) -- (ibvs.north);
\draw[flow] (ibvs.east) -- (smooth.west);
\draw[flow] (smooth.east) -- (twist.west);
\draw[flow] (twist.east) -- (robot.west);

% =================================================
% REFERENZPFADE
% =================================================
% Referenz -> Matching
\coordinate (rA) at ($(ref.east)+(0.9,0)$);
\draw[flow] (ref.east) -- (rA) |- (match.north);

% Referenz -> Bildfehler
\coordinate (rB) at ($(ref.east)+(0.45,0)$);
\draw[flow] (ref.east) -- (rB) |- (error.north);

% =================================================
% MODELLBASIERTE UNTERSTUETZUNG
% =================================================
% Prediction -> Switching
\draw[flow] (pred.north) -- (switch.south);

% Matching -> Prediction
\coordinate (mA) at ($(match.south)+(0,-0.45)$);
\draw[altflow] (match.south) -- (mA) -| (pred.west)
  node[flowlabel, pos=0.28, left] {Tracks};

% Roboterbewegung -> Prediction
\coordinate (rbA) at ($(robot.west)+(-0.9,0)$);
\draw[altflow] (robot.west) -- (rbA) |- (pred.east)
  node[flowlabel, pos=0.30, right] {Bewegung};

% =================================================
% STRECKE UND OPTISCHE RUECKKOPPLUNG
% =================================================
\draw[flow] (robot.south) -- (scene.north);

\coordinate (fbA) at ($(scene.west)+(-1.25,0)$);
\coordinate (fbB) at ($(cam.south)+(0,-0.55)$);
\draw[feedback] (scene.west) -- (fbA) |- (fbB) -| (cam.south);

% =================================================
% OPTIONALE FEINE VISUELLE HILFEN
% =================================================
% dezente Gruppierung unten
\node[font=\scriptsize, text=black!65] at (11.3,-0.95) {Vorhersage};
\node[font=\scriptsize, text=black!65] at (17.6,-0.95) {Stellgrößenaufbereitung};

\end{tikzpicture}
\end{document}